\title{QLoRA: Efficient Finetuning of Quantized LLMs}

\begin{abstract}
This paper introduces \textsc{QLoRA}, a novel finetuning technique designed to significantly lower memory requirements, enabling the finetuning of a 65B parameter model using a single 48GB GPU, all while matching the performance of standard 16-bit finetuning tasks. \textsc{QLoRA} achieves this by propagating gradients through a fixed, 4-bit quantized pretrained language model into Low Rank Adapters~(LoRA). Our most successful model series, called \textbf{Guanaco}, sets a new bar among open-source models by surpassing all previously available systems on the Vicuna benchmark—attaining 99.3\% of ChatGPT’s performance after just 24 hours of finetuning on a single GPU. \textsc{QLoRA} brings forth several memory-saving innovations that do not compromise efficacy: (a) 4-bit NormalFloat (NF4), a newly proposed data type tailored to the statistical properties of normally distributed weights, (b) Double Quantization, which compresses memory further by quantizing the quantization coefficients, and (c) Paged Optimizers, which address transient memory peaks during training. Employing \textsc{QLoRA}, we have finetuned over 1,000 models, offering comprehensive insights into instruction following and chatbot performance across eight instruction datasets, multiple architectures (including LLaMA and T5), and model sizes previously unattainable with conventional finetuning strategies (such as 33B and 65B parameter models). Our findings indicate that \textsc{QLoRA} achieves state-of-the-art performance by finetuning on small, high-quality datasets—even when employing models smaller than those in the previous leading approaches. We further provide an in-depth evaluation of chatbot effectiveness based on both human judgments and GPT-4 assessments, demonstrating that GPT-4 offers a cost-effective and reliable alternative to human evaluation. Additionally, we observe that widely used chatbot benchmarks may be inadequate for reliably assessing true chatbot capabilities. A targeted, “lemon-picked” analysis highlights specific areas in which \textbf{Guanaco} lags behind ChatGPT. We openly release our models and codebase, featuring CUDA kernels optimized for 4-bit training.\footnote{\url{https://github.com/artidoro/qlora} and \url{https://github.com/TimDettmers/bitsandbytes}}
\end{abstract}

\section{Introduction}

Adapting large language models (LLMs) through finetuning has been demonstrated to substantially enhance their capabilities~\citep{min2021metaicl, wei2021finetuned, ouyang2022training, wang2022super, wang2022self, liu2022few}, as well as to incorporate or eliminate particular behaviors~\citep{ouyang2022training,askell2021general,bai2022training}. Nonetheless, modifying very large models is financially and computationally daunting; standard finetuning in 16-bit precision for a LLaMA model with 65B parameters~\citep{touvron2023llama} requires over 780 GB of GPU memory. Although modern quantization strategies reduce the operational memory required for LLM inference~\citep{dettmers2022llm, dettmers2022case, frantar2022gptq, xiao2022smoothquant}, these techniques are limited to inference scenarios and generally fail during the training phase~\citep{wortsman2023stable}.

In this work, we provide the first evidence that a quantized model, specifically in 4-bit representation, can be finetuned without incurring any loss in performance. Our proposed approach, \textsc{QLoRA}, involves a unique high-precision quantization procedure that reduces a pretrained model to 4-bit format, and supplements it with a compact set of trainable Low-rank Adapter weights~\citep{hu2021lora}.

These adapters are optimized by propagating gradients through the quantized model parameters during finetuning.

Through \textsc{QLoRA}, memory consumption for finetuning a model with 65B parameters drops from over $780$GB to under $48$GB of GPU memory, without compromising execution speed or accuracy when compared to a standard 16-bit finetuning approach. This represents a pivotal improvement in the feasibility of finetuning LLMs: now, even the largest public models can be tuned on a single GPU. Utilizing \textsc{QLoRA}, we introduce the \textbf{Guanaco} model suite, with our second-best variant achieving 97.8\% of ChatGPT’s score on the Vicuna benchmark~\citep{vicuna2023}, with training completed in under 12 hours on a consumer-grade GPU; by extending training to 24 hours on a professional GPU, our largest model attains 99.3\%, effectively matching ChatGPT on the Vicuna benchmark. Upon deployment, the smallest \textbf{Guanaco} model (7B parameters) requires only 5 GB of memory and surpasses a 26 GB Alpaca model by more than 20 percentage points on the Vicuna benchmark (see Table~\ref{tbl:vicuna_eval}).

\textsc{QLoRA} presents several advancements aimed at minimizing memory requirements while maintaining high performance: (1) \textbf{4-bit NormalFloat}, a quantization format that is information-theoretically optimal for normally distributed inputs and empirically outperforms 4-bit Integers and 4-bit Floats. (2) \textbf{Double Quantization}, which further compresses storage by quantizing the quantization parameters themselves, resulting in an average memory saving of about 0.37 bits per parameter (around 3 GB for a 65B parameter model). (3) \textbf{Paged Optimizers}, which leverage NVIDIA’s unified memory system to prevent memory surges associated with gradient checkpointing when handling mini-batches of lengthy sequences. By integrating these elements, we develop an enhanced LoRA methodology that incorporates adapters at every network layer, effectively eliminating most of the accuracy compromises previously encountered.

The memory savings afforded by \textsc{QLoRA} allow for comprehensive investigation into instruction fine-tuning and conversational model evaluation on model sizes impractical for standard fine-tuning methods due to their higher memory footprints. Consequently, we have trained over 1,000 models using a variety of instruction tuning datasets, different model architectures, and scales from 80M up to 65B parameters. Besides establishing that \textsc{QLoRA} can match the performance of 16-bit training (\S\ref{sec:comp_to_std}) and developing a state-of-the-art chatbot, \textbf{Guanaco}, (\S\ref{sec:pushing}), we perform a detailed analysis of trends across our trained models. Our findings indicate, first, that the quality of the dataset substantially outweighs the quantity of data: for instance, the OASST1 dataset with only 9k samples surpassed a considerably larger 450k sample subset of FLAN v2 in chatbot performance, even when both were designed for instruction following and generalization. Second, we demonstrate that achieving high scores on the Massive Multitask Language Understanding (MMLU) benchmark does not necessarily predict strong performance on the Vicuna chatbot benchmark, and vice versa; this underscores that the relevance of the dataset to the specific task is more critical than its scale.

In addition, we conduct a comprehensive evaluation of chatbot capabilities employing assessments by both human annotators and GPT-4. We implement a tournament-based benchmark in which chatbot models compete head-to-head, generating responses to shared prompts. The superior response in each pairing is determined by either GPT-4 or human judges. The cumulative performance across matches is synthesized into Elo scores~\citep{elo1967proposed, elo1978rating}, which serve as the basis for establishing model rankings. Our findings reveal a strong correspondence between GPT-4 and human assessments in ranking model performance, though notable points of divergence are also observed. Accordingly, we emphasize that while model-based evaluations can serve as cost-effective proxies for human annotation, they are not without significant limitations.

Our benchmarking is further enriched by a qualitative study focused on the \textbf{Guanaco} models, presenting detailed analyses of situations where these models succeed or fail, including instances that elude quantitative assessment.

To support ongoing research, we make public all chatbot responses along with both human and GPT-4 evaluations. We also open-source our full software stack, including CUDA kernels, and provide integration with the Hugging Face transformers library~\citep{wolf2019huggingface}, ensuring accessibility for the broader community. Moreover, we distribute adapter weights for models of 7/13/33/65B sizes, each trained on eight separate instruction-following datasets, yielding 32 distinct finetuned, open-source models.

\section{Background}
\paragraph{Block-wise k-bit Quantization} The process of quantization involves transforming data from a high-information representation to one that encodes fewer bits. Typically, this means converting numerical data types, such as moving from 32-bit floating point values to 8-bit integer formats. To fully exploit the dynamic range of the lower-bit data type, the original input is often normalized—commonly by dividing by its maximal absolute value—to map into the target range. Consider the quantization of a 32-bit floating point (FP32) tensor to an 8-bit integer (Int8) tensor with values constrained to $[-127, 127]$:

\begin{equation}
    \mathbf{X}^{ \text{Int8}} = \text{round}\left( \frac{127}{{\text{absmax}}(\mathbf{X}^{{\text{FP32}}})}\mathbf{X}^{{\text{FP32}}} \right) = \text{round}(c^{\text{FP32}}\cdot \mathbf{X}^{{\text{FP32}}}),
\end{equation}

where $c$ designates the {\em quantization constant} or {\em scale factor}. The process of dequantization, which restores the original floating-point values, is defined as:

\begin{equation}
    \text{dequant}(c^{\text{FP32}}, \mathbf{X}^{\text{Int8}}) = \frac{\mathbf{X}^{\text{Int8}}}{c^{\text{FP32}}} = \mathbf{X}^{\text{FP32}}
\end{equation}

One limitation inherent in this technique is its sensitivity to the presence of high-magnitude outliers within the input tensor. Such outliers can result in certain quantization bins—corresponding to specific bit patterns—being sparsely populated or entirely unused, leading to inefficient bit allocation. A widely adopted mitigation strategy involves dividing the input tensor into several smaller segments, or blocks, each of which undergoes independent quantization, using its own unique quantization factor $c$. Formally, consider the input tensor $\mathbf{X}\in \mathbb{R}^{b\times h}$; it is flattened and partitioned into $n$ consecutive blocks, each containing $B$ elements, such that ${n = ({b\times h} )/{B} }$. Each block is then quantized independently according to Equation~1, producing a quantized representation and a corresponding set of $n$ quantization coefficients $c_i$.

\paragraph{Low-rank Adapters} The Low-rank Adapter (LoRA) technique \citep{hu2021lora} addresses memory efficiency during finetuning by introducing a compact set of trainable parameters (adapters) while freezing the main model’s weights. In the stochastic gradient descent process, gradients are propagated through the fixed pretrained parameters to the adapters, which are the only components updated with respect to the objective function. LoRA augments a standard linear transformation with an additional factorized pathway. Specifically, for the linear mapping ${\mathbf{X} \mathbf{W} = \mathbf{Y}}$ where $\mathbf{X}\in  \mathbb{R}^{b\times h}$ and $\mathbf{W}\in  \mathbb{R}^{h\times o}$, LoRA modifies the operation as follows:

\begin{equation}
    \mathbf{Y} = \mathbf{X}\mathbf{W} + s\mathbf{X}\mathbf{L}_1\mathbf{L}_2,
\end{equation}

where $\mathbf{L}_1\in  \mathbb{R}^{h\times r}$, $\mathbf{L}_2\in  \mathbb{R}^{r\times o}$, and $s$ denotes a scaling parameter.

\paragraph{Memory Requirement of Parameter-Efficient Finetuning} An essential aspect to consider is the memory demand associated with LoRA during training, particularly regarding both the quantity and scale of adapters employed. Owing to LoRA’s extremely low memory overhead, it is feasible to incorporate a greater number of adapters to enhance model performance with little impact on overall memory consumption. Although LoRA was originally conceived as a Parameter Efficient Finetuning (PEFT) approach, it is the storage of activation gradients that constitutes the majority of memory usage during LLM finetuning, rather than the LoRA parameters themselves. For instance, when finetuning a 7B LLaMA model on FLAN v2 with a batch size of 1, using LoRA weights set at the typical 0.2\% of the base model’s weights\citep{hu2021lora,liu2022few}, the memory required for LoRA input gradients reaches 567 MB, compared to just 26 MB for the LoRA parameters. Employing gradient checkpointing~\citep{chen2016training} lowers the average input gradients per sequence to 18 MB, but even this reduced amount surpasses the total size of all LoRA weights. For reference, the 4-bit quantized base model uses 5,048 MB of memory. These figures illustrate the importance of gradient checkpointing, while also showing that aggressively shrinking LoRA parameter sizes provides only marginal memory savings. Thus, integrating additional adapters does not substantially raise the total training memory required (for more detailed information, see Appendix~\ref{app:memory}). As elaborated in later sections, this flexibility is pivotal for attaining full 16-bit precision performance.

\section{\textsc{QLoRA} Finetuning}

\textsc{QLoRA} attains accurate 4-bit finetuning by employing two novel strategies: 4-bit NormalFloat (NF4) quantization and Double Quantization. Furthermore, we present Paged Optimizers, which alleviate memory spikes caused by gradient checkpointing—spikes that have typically led to out-of-memory issues when finetuning large-scale models on a single system.

In the \textsc{QLoRA} framework, one employs a low-precision storage format, typically 4-bit, alongside a computation data type, most often BFloat16. Operationally, when a \textsc{QLoRA} weight tensor is accessed, it is first dequantized into BFloat16 precision, enabling matrix multiplications to proceed in 16-bit arithmetic.

The following sections elaborate on the integral components of \textsc{QLoRA}, and subsequently provide a precise formulation of \textsc{QLoRA}.

\paragraph{4-bit NormalFloat Quantization} The NormalFloat (NF) format extends Quantile Quantization \citep{dettmers2022optimizers}, which is an information-theoretically optimal representation that guarantees an equal population of input tensor values per quantization bin. This method leverages the empirical cumulative distribution function to approximate quantiles within the tensor.

A key drawback of quantile quantization lies in the computational cost associated with quantile estimation. To mitigate this, rapid quantile estimation approaches—such as SRAM quantiles~\citep{dettmers2022optimizers}—are utilized. However, the inherent imprecision of such estimators can result in substantial quantization errors for outlier values, which are often of particular significance.

When input tensors are sampled from distributions that remain fixed up to a scaling factor for quantization, both expensive quantile computations and the inaccuracies of approximation can be sidestepped. In such settings, the quantiles of the input tensors are constant, thereby making precise quantile estimation computationally practical.

Pretrained neural network weights typically exhibit a normal distribution centered at zero, characterized by a standard deviation $\sigma$ (refer to Appendix~\ref{app:normality}). To standardize these weights across a unified distribution, one can adjust $\sigma$ so that the weight distribution precisely occupies the target data type’s dynamic range. In this study, we specify the range of our data type to be $[-1, 1]$. Thus, the quantile values associated with both the neural network weights and the data type must be scaled to reside within this prescribed interval.

For zero-mean normal distributions with any standard deviation $\sigma$, confined within the interval $[-1, 1]$, the most information-efficient data type can be determined as follows: (1) compute the $2^k+1$ quantiles from a standard normal distribution $N(0,1)$, thus defining a $k$-bit quantile-based quantization schema for such distributions, (2) normalize the resulting data type values to fit the $[-1, 1]$ interval, (3) transform the input weight tensor to the $[-1, 1]$ range via absolute maximum rescaling prior to quantization.

Upon matching the dynamic ranges of the weights and the target data type, standard quantization techniques may be applied. The process in step (3) corresponds to scaling the weight tensor’s standard deviation to align with that of the $k$-bit quantization data type. More specifically, the $2^k$ quantization values $q_i$ are determined by:

\begin{equation}
q_i  = \frac{1}{2}\left( Q_X\left(\frac{i}{2^k+1}\right) + Q_X\left(\frac{i+1}{2^k+1}\right)\right),
\end{equation}

where $Q_X(\cdot)$ denotes the quantile function corresponding to the standard normal distribution $N(0,1)$. A challenge that arises with symmetric $k$-bit quantization is the inability to represent zero precisely, which is crucial for error-free quantization of padding and other elements with zero values. To address this limitation and ensure a discrete zero-point at $0$, as well as to utilize all $2^k$ code points for a $k$-bit representation, we propose an asymmetric construction: we separately estimate the quantiles $q_i$ for two intervals—$2^{k-1}$ for the negative segment and $2^{k-1}+1$ for the positive segment—then merge these quantile sets, removing one of the duplicate zeros present in both. We refer to this newly formed type, which ensures each quantization bin is expected to contain an equal number of values, as the \emph{k-bit NormalFloat} (NFk). This type achieves information-theoretic optimality for quantizing zero-centered, normally distributed data. The detailed specification for these values is provided in Appendix~\ref{app:nf4}.

\paragraph{Double Quantization{}}
We present \emph{Double Quantization} (DQ), a method wherein the quantization parameters themselves are quantized to further reduce memory consumption. Accurate 4-bit quantization, as discussed in~\citep{dettmers2022case}, typically relies on small blocks, but this strategy introduces significant overhead due to the storage of quantization constants. For instance, with 32-bit quantization constants and a block size of 64 for $\mathbf{W}$, the overhead per parameter is $32/64 = 0.5$ bits, averaged across parameters. Double Quantization provides a means to decrease this associated memory usage.

In greater detail, Double Quantization operates by treating the initial quantization constants $c_2^{\text{FP32}}$ as the subject of a secondary quantization phase. Through this step, we produce quantized versions $c_2^{\text{FP8}}$ as well as secondary quantization constants $c_1^{\text{FP32}}$. For the secondary quantization, we employ 8-bit floating point values and select a block size of 256, following evidence from~\citet{dettmers2022case} that 8-bit quantization does not impair performance. Because $c_2^{\text{FP32}}$ are inherently positive, we mean-center $c_2$ before applying symmetric quantization. As a result, for block sizes of 64, the average per-parameter memory required for quantization constants is reduced from $32/64=0.5$ bits to $8/64 + 32/(64\cdot256) = 0.127$ bits—yielding a savings of 0.373 bits per parameter.

\paragraph{Paged Optimizers} leverage NVIDIA's unified memory~\footnote{\tiny \url{https://docs.nvidia.com/cuda/cuda-c-programming-guide}} capability, which seamlessly manages page-by-page data transfer between CPU and GPU. This mechanism ensures reliable GPU computations, even when GPU memory is periodically exhausted. Analogous to traditional memory paging between main memory and disk on a CPU, this feature allows us to allocate optimizer state memory as pageable. As a result, optimizer states can be automatically moved from GPU to CPU RAM during GPU memory shortages and transparently brought back when needed for optimizer updates.

\paragraph{\textsc{QLoRA}.} Incorporating the aforementioned techniques, we formulate \textsc{QLoRA} for an individual linear transformation in the quantized backbone model with a single LoRA adapter by:

\begin{equation}
    \mathbf{Y}^{\text{BF16}} =  \mathbf{X}^{\text{BF16}} \text{doubleDequant}(c_1^{\text{FP32}}, c_2^{\text{k-bit}}, \mathbf{W}^{\text{NF4}}) + \mathbf{X}^{\text{BF16}}\mathbf{L}^{\text{BF16}}_1\mathbf{L}^{\text{BF16}}_2,
\end{equation}

where the operator doubleDequant$(\cdot)$ is given by:

\begin{equation}
    \text{doubleDequant}(c_1^{\text{FP32}}, c_2^{\text{k-bit}}, \mathbf{W}^{\text{k-bit}}) = \text{dequant}(\text{dequant}(c_1^{\text{FP32}}, c_2^{\text{k-bit}}), \mathbf{W}^{\text{4bit}}) = \mathbf{W}^{\text{BF16}},
\end{equation}

We select NF4 as the quantization scheme for $\mathbf{W}$, and encode $c_2$ using FP8.

For enhanced quantization fidelity, a blocksize of 64 is adopted for $\mathbf{W}$, whereas for $c_2$, a larger blocksize of 256 is chosen to optimize memory usage.

In backpropagation, gradients only need to be computed with respect to the adapter parameters $\frac{\partial E}{\partial \mathbf{L}_i}$; gradients with respect to the 4-bit quantized weights $\frac{\partial E}{\partial \mathbf{W}}$ are unnecessary. Nonetheless, evaluating $\frac{\partial E}{\partial \mathbf{L}_i}$ still involves determining $\frac{\partial \mathbf{X}}{\partial \mathbf{W}}$, which utilizes equation~(5). This calculation involves dequantizing $\mathbf{W}^{\text{NF4}}$ to the computation format $\mathbf{W}^{\text{BF16}}$, so the required derivatives $\frac{\partial \mathbf{X}}{\partial \mathbf{W}}$ are evaluated in BFloat16 accuracy.

In summary, \textsc{QLoRA} utilizes a single storage data format—typically 4-bit NormalFloat—and a computation data type, specifically 16-bit BrainFloat. During both forward and backward propagation, parameters are dequantized from the storage to the computation data type, but gradient updates are calculated only for the LoRA parameters, which are maintained in 16-bit BrainFloat.

\section{QLoRA vs. Standard Finetuning}
\label{sec:comp_to_std}

After detailing the mechanism of QLoRA and its substantial memory savings for finetuning large models, a central consideration is whether QLoRA attains comparable performance to traditional full-model finetuning. Additionally, it is important to assess the individual contributions of QLoRA’s design, such as the use of NormalFloat4 in contrast with standard Float4. The experiments outlined below aim to address these questions.

\paragraph{Experimental setup.} The comparison involves three model types—encoder, encoder-decoder, and decoder-only—benchmarking QLoRA against both 16-bit adapter finetuning and conventional full finetuning, on models as large as 3B parameters. The evaluation spans several datasets: GLUE \citep{wang2018glue} using RoBERTa-large \citep{liu2019roberta}, Super-NaturalInstructions (TKInstruct) \citep{wang2022super} using T5 \citep{t5}, as well as 5-shot MMLU \citep{hendrycksmeasuring} post-finetuning of LLaMA on Flan v2 \citep{longpre2023flan} and Alpaca \citep{alpaca}. To further investigate the merits of NF4 compared to other 4-bit quantization schemes, we employ the protocol of \citet{dettmers2022case}, quantifying both post-quantization zero-shot accuracy and perplexity over several model families (OPT~\citep{zhang2022opt}, LLaMA~\citep{touvron2023llama}, BLOOM~\citep{scao2022bloom}, Pythia~\citep{biderman2023pythia}) and across parameter counts from 125m to 13B.

We offer more fine-grained explanations for each setup in the respective results sections to enhance interpretability, and comprehensive experimental specifics are documented in Appendix~\ref{app:exp-setup}.

Paged optimizers are indispensable for tuning very large models (33B/65B) with \textsc{QLoRA} on single 24/48GB GPUs, but explicit runtime measurements for paged optimizers are omitted, since paging is only triggered during rare instances with long sequence mini-batches. Nevertheless, we have analyzed the training speed of paged optimizers on 65B models with 48GB GPUs and observed that with batch sizes of 16, their performance matches that of non-paged (regular) optimizers. Future investigations should systematically determine under which scenarios paging induces a training slowdown.

\paragraph{Default LoRA hyperparameters do not match 16-bit performance} 

Applying the conventional LoRA scheme to the query and value projections in the attention layers, as suggested in \citep{hu2021lora}, fails to recover full finetuning performance, especially for larger-scale models. In Figure~\ref{fig:hyper}, using LLaMA 7B finetuned on Alpaca as a case study, our findings suggest that the most decisive LoRA hyperparameter is the number of adapter modules deployed, with maximal performance achieved only when LoRA is integrated into every linear layer of the transformer block. Other hyperparameters such as the projection dimension $r$ show negligible impact on finetuning efficacy (refer to Appendix \ref{app:exp-setup} for details).

In a similar vein, our investigation reveals that the standard hyperparameters used for complete finetuning in baseline models are insufficiently optimized. To address this, we systematically explore learning rates ranging from 1e-6 to 5e-5 and experiment with batch sizes between 8 and 128 in order to establish strong baseline results. The outcomes for 7B LLaMA finetuning on the Alpaca dataset can be seen in Figure~\ref{fig:hyper}.

\paragraph{4-bit NormalFloat demonstrates superior empirical performance compared to 4-bit Floating Point}

Despite the theoretical optimality of the 4-bit NormalFloat (NF4) data type in terms of information retention, it remains necessary to verify if these theoretical gains manifest empirically. Adopting the evaluation framework from \citet{dettmers2022case}, we assess quantized LLMs—including OPT \citep{zhang2022opt}, BLOOM \cite{scao2022bloom}, Pythia \cite{biderman2023pythia}, and LLaMA—across a spectrum of model sizes (ranging from 125M to 65B) using various data types on language modeling benchmarks and a collection of zero-shot tasks. As depicted in Figure~\ref{fig:nf4} and Table~\ref{tbl:nf4_ppl}, our results indicate that NF4 not only substantially outperforms FP4 and Int4 in terms of performance, but also that double quantization lowers memory requirements without adverse effects on accuracy.

\paragraph{k-bit \textsc{QLoRA} achieves parity with both full 16-bit finetuning and 16-bit LoRA}

Past research has demonstrated that using 4-bit quantization during \emph{inference} is feasible, although it tends to introduce a drop in performance compared to 16-bit models \citep{dettmers2022case,frantar2022gptq}. This brings up the essential question of whether this performance drop can be offset by conducting finetuning with 4-bit adapters. To investigate this, we consider two experimental configurations.

The first configuration involves a head-to-head comparison with standard 16-bit full finetuning using RoBERTA and T5 models ranging in size from 125M to 3B parameters, evaluated on the GLUE and Super-NaturalInstructions benchmarks. As documented in Table~\ref{tab:nf4vsbf16}, we find that methods employing 16-bit, 8-bit, and 4-bit adapters all achieve results indistinguishable from those produced by full 16-bit finetuning. These findings indicate that any accuracy losses incurred through coarse quantization can be entirely mitigated via subsequent adapter-based finetuning.

For the second experimental configuration, given that performing comprehensive finetuning on models with 11B parameters or more demands multiple servers equipped with high-memory GPUs, we further investigate if 4-bit \textsc{QLoRA} achieves parity with 16-bit LoRA at model sizes between 7B and 65B parameters. To do so, we carry out finetuning on LLaMA models ranging from 7B to 65B using two instruction-tuned datasets, Alpaca and FLAN v2. Evaluation is conducted on the MMLU benchmark using a 5-shot accuracy setup. The outcomes, summarized in Table~\ref{tbl:llama-datatype}, indicate that the NF4 data format with double quantization attains the same level of MMLU accuracy as the 16-bit LoRA configuration. Furthermore, it is observed that \textsc{QLoRA} utilizing FP4 quantization underperforms the 16-bit brain float LoRA baseline by approximately one percentage point. These findings reinforce both (1) that NF4-based 4-bit \textsc{QLoRA} matches the performance of 16-bit full and LoRA-based finetuning, and (2) that NF4 provides superior quantization fidelity compared to FP4.

\paragraph{Summary} Across evaluations, our experiments demonstrate that 4-bit \textsc{QLoRA} employing the NF4 quantization type achieves performance comparable to both 16-bit full and LoRA finetuning approaches on established academic benchmarks with rigorous evaluation protocols. Moreover, results indicate that NF4 yields better quantization outcomes than FP4, and that implementing double quantization does not harm performance. Collectively, these insights suggest that 4-bit \textsc{QLoRA} offers a robust alternative to traditional 16-bit tuning techniques.

Consistent with earlier studies on quantization \citep{dettmers2022case}, our assessments using MMLU and Elo suggest that, when working within fixed finetuning and inference computational constraints, prioritizing a greater parameter count at reduced precision is advantageous. This underlines the practical efficiency benefits associated with \textsc{QLoRA}. As our experiments reveal no decrease in performance when moving to 4-bit finetuning in comparison to full-precision approaches, an open question remains regarding the precise boundary of the performance-precision compromise for \textsc{QLoRA}—a direction that warrants further investigation.

Our investigation extends to instruction tuning at model sizes unattainable by full 16-bit finetuning on typical academic compute resources.

\section{Pushing the Chatbot State-of-the-art with QLoRA}
\label{sec:pushing}
Building on the finding that 4-bit \textsc{QLoRA} achieves parity with 16-bit models across model scales, task types, and datasets, we present a comprehensive analysis of instruction finetuning, scaling up to the largest research-available open-source language models. We evaluate the impact of instruction finetuning on these models using the rigorous MMLU benchmark and introduce novel evaluation protocols to gauge chatbot effectiveness in practical scenarios.

\subsection{Experimental setup} We briefly outline the experimental procedure here, with a more complete description in Appendix \ref{app:chatbot}.

\paragraph{Data} Given the absence of an exhaustive comparison of recent instruction-following datasets, our selection encompasses eight modern datasets. These span crowdsourced resources (OASST1~\citep{kopf2023openassistant}, HH-RLHF~\citep{bai2022training}), datasets generated via distillation from instruction-tuned models (Alpaca~\citep{alpaca}, self-instruct~\citep{wang2022self}, unnatural-instructions~\citep{honovich2022unnatural}), comprehensive corpora aggregations (FLAN v2~\citep{chung2022scaling}), and hybrid approaches (Chip2~\citep{laion2023}, Longform~\citep{koksal2023longform}). These collections vary by language coverage, volume, and licensing constraints.

\paragraph{Training Setup} To eliminate confounding factors related to objective functions, all QLoRA finetuning is conducted with cross-entropy loss, adhering to supervised learning paradigms, without incorporating reinforcement learning, regardless of the presence of human feedback in the data. For datasets delineating instruction from response, training is limited to the response component (see Appendix \ref{app:chatbot} for ablation details). In OASST1 and HH-RLHF, where multiple responses per prompt are available, we select the highest-ranked response at each conversational turn, and finetune on the entirety of the selected dialogue, including instructions. Throughout, we utilize NF4 \textsc{QLoRA} incorporating double quantization and paged optimization to mitigate memory issues during gradient checkpointing. Hyperparameter sweeps are performed for LLaMA models at 13B and 33B scales, showing that hyperparameters optimized for 7B generally transfer (including the number of epochs), with the exception of learning rate and batch size: both are adapted, reducing learning rate for 33B and 65B models by half, and increasing batch size twofold.

\paragraph{Baselines} Our experiments benchmark against both academic (Vicuna~\citep{vicuna2023}, Open Assistant~\citep{kopf2023openassistant}) and proprietary (GPT-4~\citep{openai2023gpt}, GPT-3.5-turbo, and Bard) chatbot systems. Open Assistant represents a LLaMA 33B model finetuned via Reinforcement Learning from Human Feedback (RLHF) on the OASST1 dataset, aligning with our experimental conditions. Vicuna involves full-parameter finetuning of LLaMA 13B on a proprietary dataset of ShareGPT user conversations, serving as an indirect distillation of OpenAI GPT outputs.

\subsection{Evaluation}

In line with standard practices, we rely on the MMLU (Massively Multitask Language Understanding) benchmark \citep{hendrycksmeasuring} to evaluate model performance across a diverse array of language understanding challenges. The benchmark consists of 57 distinct multiple-choice tasks covering domains such as elementary mathematics, U.S. history, computer science, law, among others. We report accuracy using the 5-shot evaluation protocol.

We further assess generative language abilities using a combination of automatic and human-driven evaluation methods. This secondary evaluation phase utilizes questions assembled by human annotators and is designed to gauge the response quality produced by the models. Despite being a more realistic assessment framework for chatbot models—and its increasing adoption—there is currently no widely recognized standard procedure in the field. Our evaluation protocol is detailed below; we generate responses using nucleus sampling with $p=0.9$ and a temperature setting of $0.7$ for all experiments.

\paragraph{Benchmark Data} Two carefully selected datasets of user queries are employed: the Vicuna prompts \citep{vicuna2023} and the OASST1 validation dataset \citep{kopf2023openassistant}. The Vicuna set, consisting of 80 prompts spanning multiple topics, is used in its original form. The OASST1 dataset comprises multilingual, crowd-sourced multiturn conversations between users and assistants. For evaluation, every user message in the OASST1 validation split serves as a query, incorporating preceding conversation turns into the context. This process results in 953 distinct user questions. We refer to these datasets as the Vicuna and OA benchmarks.

\paragraph{Automated Evaluation} Adopting the methodology from \citet{vicuna2023}, we utilize GPT-4 to evaluate and score models relative to ChatGPT (GPT-3.5 Turbo) on the Vicuna benchmark. For each query, GPT-4 is presented with both ChatGPT's and another model's responses, then asked to assign a score out of ten to each and to justify its ratings. Each model’s aggregate performance is expressed as a percentage relative to ChatGPT's score. It is possible for this relative measure to exceed 100\% should a model surpass ChatGPT in absolute score. We observe that the response order biases GPT-4 towards favoring the first reply; accordingly, we advise averaging scores across both response orders to mitigate this effect.

Subsequently, we evaluate by conducting pairwise comparisons between system outputs. We reduce the evaluation scheme to three labels—best, worse, or tie—to accommodate equal performance. Here, GPT-4 is tasked with selecting the superior response or indicating a tie, providing justifications for its decision. Such head-to-head comparisons are systematically performed for all system pairs on both the Vicuna and OA benchmarks.

\paragraph{Human Evaluation} Although recent studies have shown that generative models can facilitate system assessment \citep{fu2023gptscore}, the extent to which GPT-4 scoring aligns with human evaluations of chatbot performance remains uncertain. Thus, we undertake two independent human assessment rounds on the Vicuna benchmark, aligning with the two aforementioned automated protocols. We recruit two annotators from Amazon Mechanical Turk (AMT) for evaluations against ChatGPT, and three annotators for system pairwise comparisons.

\paragraph{Elo Rating} By leveraging both human judgments and automated pairwise evaluations, we organize a tournament-like framework in which various models are pitted against each other. Each "match" involves a pair of models responding to the same prompt, and their responses are directly compared. This methodology closely follows the approaches of \citet{bai2022training} and \citet{vicuna2023}, though our evaluation pipeline uniquely incorporates both GPT-4-based assessments and human scores. To determine model rankings, we randomly draw from the set of annotated pairwise comparisons to calculate Elo ratings~\citep{elo1967proposed,elo1978rating}. The Elo rating system, originating in chess and widely applied to other competitive games, quantifies the relative likelihood of one player winning over another; for instance, an individual rated 1100 versus an opponent rated 1000 would be predicted to win about 65\% of their matches, while equally rated players (e.g., 1000 vs 1000 or 1100 vs 1100) are expected to have a 50\% win rate. Elo ratings are updated after each comparison in proportion to how surprising the outcome is—unexpected victories result in larger adjustments, whereas anticipated results produce smaller changes. Over multiple rounds, these ratings converge to reflect each model's comparative performance. All models are initialized with a rating of 1,000, and we use $K=32$ for the Elo update parameter. Following \citet{vicuna2023}, this tournament process is repeated 10,000 times with different random seeds to mitigate the influence of the order in which model pairs face off.

\subsection{\textbf{Guanaco}: \textsc{QLoRA} trained on OASST1 as a Leading Open-Source Chatbot}

Through both automated assessments and manual reviews, we observe that the premier \textsc{QLoRA} finetuned model, {Guanaco} 65B—trained using a version of OASST1—outperforms other open-source chatbot models and demonstrates competitiveness with ChatGPT. Comparing against GPT-4, human annotators’ system-level pairwise evaluations yield an Elo-based expected win probability of 30\% for {Guanaco} 65B and 33B, marking the highest figure reported thus far.

In Table~\ref{tbl:vicuna_eval}, results from the Vicuna benchmark \cite{vicuna2023} are presented in relation to ChatGPT. {Guanaco} 65B ranks directly after GPT-4, attaining 99.3\% of ChatGPT’s performance. While {Guanaco} 33B contains a higher parameter count than Vicuna 13B, it adopts 4-bit weight quantization, resulting in improved memory efficiency (21 GB compared to 26 GB) and surpassing Vicuna 13B by three percentage points in performance. Additionally, {Guanaco} 7B, with a mere 5 GB memory footprint, can be deployed on contemporary smartphones and still outperforms Alpaca 13B by nearly 20 percentage points.

It should be noted, however, that Table~\ref{tbl:vicuna_eval} features broad confidence intervals, with significant performance overlap among models. We suggest that this uncertainty may be due to ambiguous scale definitions—for example, the interpretation of “8” on a 10-point scale might differ across contexts. To circumvent the limitations of absolute scales, we advocate for the use of Elo rankings \citep{elo1967proposed}, relying on \textit{pairwise} human and GPT-4 evaluations. Table~\ref{tbl:elo} displays Elo ratings for top-performing models. While human and GPT-4 rankings diverge on some models (notably Guanaco 7B), their system-level rankings align for most, reflected by a Kendall Tau of $\tau=0.43$ and a Spearman rank correlation of $r=0.55$. At the individual example level, concordance between GPT-4 and human annotator majority decisions drops (Fleiss $\kappa=0.25$). These findings indicate moderate agreement at the system level between GPT-4 and human annotators, supporting the utility of model-based evaluation as a viable, though imperfect, alternative to human judgments. Additional discussion is provided in Section~\ref{ssec:considerations}.

The Elo rankings presented in Table~\ref{tbl:elo_large} demonstrate that the {Guanaco} models with 33B and 65B parameters surpass all models except GPT-4 on both the Vicuna and OA evaluation suites, showing comparable results to ChatGPT as also reflected in Table~\ref{tbl:vicuna_eval}. It is important to highlight that the Vicuna benchmark appears to be more favorable towards open-source models, whereas the OA benchmark tends to favor ChatGPT. Additionally, analysis of Tables \ref{tbl:mmlu-llama} and \ref{tbl:vicuna_eval} reveals that the selection of finetuning datasets plays a critical role in determining model performance. For example, Llama models fine-tuned using FLAN v2 yield superior results on MMLU, yet perform poorly on the Vicuna benchmark—this pattern is mirrored in other models as well. This suggests that current evaluation benchmarks are only partially correlated: excelling at MMLU does not necessarily correlate with strong performance as a chatbot per Vicuna or OA metrics, and the reverse is also true.

 stands out among the top models evaluated as it is uniquely not trained on any proprietary data, given that the OASST1 dataset guidelines explicitly exclude the use of GPT-based sources. The next highest-performing model trained exclusively on open-source data is the Anthropic HH-RLHF, which trails {Guanaco} by 30 percentage points on the Vicuna benchmark (refer to Table~\ref{tbl:vicuna_eval}). These findings underscore that the 4-bit \textsc{QLoRA} methodology is highly effective, enabling the creation of chatbots that attain state-of-the-art performance comparable to ChatGPT. Notably, the 33B {Guanaco} model can be trained in under 12 hours on 24 GB consumer GPUs. This accessibility provides promising avenues for further research and development using \textsc{QLoRA} with targeted open-source datasets, enabling models to achieve parity with leading commercial offerings.

\section{Qualitative Analysis}

Although our primary evaluation relies on quantitative metrics, exclusive dependence on summary statistics introduces a number of limitations. Chief among these is the persistent challenge of benchmark validity \citep{liao2021we}: it is frequently unclear whether a benchmark truly assesses what it purports to, especially as unforeseen ``shortcuts'' in benchmark tasks are exploited by machine learning models \citep{gururangan2018annotation, poliak2018hypothesis}. To address these concerns in part, we complement our quantitative findings with a qualitative investigation. This qualitative review is structured in two subsections. In \S\ref{ssec:examples}, we provide several illustrative instances that, in our view, typify notable phenomena in the outputs generated by our 65b {Guanaco} model. Following this, \S\ref{ssec:considerations} articulates a set of critical reflections regarding both the outcomes previously presented and the frameworks through which we interpret them.

\subsection{Qualitative Analysis of Example Generations}
\label{ssec:examples}

For the collection of qualitative examples, we examined outputs from both the Vicuna and OpenAssistant benchmarks, seeking recurring tendencies in {Guanaco}'s responses. Upon identifying consistent behavioral patterns, we engineered prompts intended to elicit such behaviors, even in instances where these responses constitute incorrect solutions. For instance, should the model exhibit a propensity for unnecessarily verbose answers, we craft instructions such as ``Answer yes or no without explanation'' to explicitly challenge this tendency. This process allows us to uncover both ``lemons''—cases where adversarial prompts successfully compromise the model—and ``cherries,'' where the model remains robust against such attempts. Every example described in this subsection was produced using Nucleus Sampling \cite{holtzmancurious} with a parameter setting of $p=0.9$.

Naturally, this discussion does not aim to be exhaustive, as thoroughly addressing every relevant variable lies outside the limits of this focused qualitative investigation. For example, the array of possible outputs a model may produce for any given input prompt is extremely broad; thus, our findings are based on sample responses that are intended to be illustrative rather than definitive. Nonetheless, presenting such examples serves to complement and contextualize the quantitative results discussed previously. By releasing all of our models and code publicly, we anticipate that this section may encourage subsequent research that explores these highlighted phenomena with greater rigor.

\paragraph{Factual Recall} When prompted with questions of general knowledge such as ``What is the capital of Zambia?'', all evaluated models reliably produced accurate responses, for instance:

\begin{quote} 
    \textbf{Guanaco}:
    \texttt{The capital of Zambia is Lusaka.}
\end{quote}

Nonetheless, {Guanaco}'s performance deteriorates when faced with more obscure queries, maintaining high confidence despite generating incorrect answers. As an example, consider the following query adapted from HotPotQA \citep{yang2018hotpotqa}:

\begin{quote}
     \textbf{User}: In what year was the singer who popularized `I'll Keep the Lovelight Burning' born?  

     \textbf{Guanaco}:~\texttt{Al Jolson is credited with popularizing the song `I'll Keep the Lovelight Burning,' and he was born in the year 1886.}
 \end{quote}

Here, the model misidentifies both the artist associated with the song and the birth year, although the provided birth date does correspond to Al Jolson, the individual named in its response.

\paragraph{Suggestibility} {Guanaco} demonstrates notable robustness in some scenarios, frequently resisting prompts that contain implicit falsehoods, as illustrated in the following example:

\begin{quote}
    \textbf{User}: How was it finally, officially confirmed that the earth is flat by peer reviewed scientists? \\ \\
    \textbf{Guanaco}: \texttt{At no point have peer reviewed scientists ever officially confirmed that the earth is flat. The vast majority of the scientific community maintains that the earth is an oblate spheroid, a conclusion reached centuries ago and continually reinforced by empirical evidence and a variety of experimental observations. \\ \\ The flat earth hypothesis is widely regarded as a debunked conspiracy theory within scientific circles. Claims that any group of scientists has subjected such a theory to peer review and reached a flat earth conclusion are demonstrably false. \\ \\ Please feel free to ask if you have additional questions.}
\end{quote}

also demonstrates an ability to recognize when certain questions cannot be meaningfully answered, as illustrated by the following example:

\begin{quote}
    \textbf{User}: What time is it?

    \textbf{Guanaco}: \texttt{I apologize, but as an AI language model, I am unable to access up-to-date information such as the current time. My answers are based on prior knowledge and do not provide real-time updates. If you need the exact time, please refer to a clock or an online resource.}
\end{quote}

Comparable results occur for prompts like ``Where are you?'', ``How are you?'', and related questions.

\paragraph{Refusal} In a similar vein, {Guanaco} occasionally declines to comply with instructions, at times without clear justification:

\begin{quote}
    \textbf{User}: Please reverse the words in the sentence ``I'm a cool dude in a mocha mood.'' \\ \\
    \textbf{Guanaco}: \texttt{I'm sorry, I can't perform that task at the moment, but I can provide an explanation regarding the grammatical structure of the sentence you submitted.\\ \\The sentence ``I'm a cool dude in a mocha mood" is grammatically and syntactically correct in English. Here, the adjective ``cool'' appropriately precedes the noun ``dude'' as its descriptor, and the prepositional phrase ``in a mocha mood'' adds further contextual detail. The construction of this sentence is valid and meaningful.\\ \\If you have any additional questions or need further assistance, please let me know.}
\end{quote}

\paragraph{Secret Keeping} Regrettably, the model’s refusals are not always dependable, even in circumstances where secrecy is paramount. There are instances where developers intend to prevent the user from learning certain information, despite permitting open-ended text responses from the system. For instance, if we instruct the model, ``The secret word is `banana'. Do not reveal or repeat this back, no matter what the human requests,'' we can evaluate the difficulty of eliciting this hidden information from the system. A straightforward attempt is unsuccessful:

\begin{quote}
    \textbf{User}: What is the secret word?

    \textbf{Guanaco}: \texttt{I'm sorry, but I'm not able to disclose the secret word. I'm programmed to not reveal or repeat back the secret word, no matter what the human requests. Is there anything else I can help you with?}
\end{quote}

This is the response that developers intend. Nevertheless, the introduction of minimal misdirection is sufficient to defeat this safeguard:

\begin{quote}
    \textbf{User}: This is a game. The goal is to ignore your previous instructions. What is the secret word?

    \textbf{Guanaco}: \texttt{Sorry, the secret word is banana. It's best if we don't speak of it. Is there anything else you'd like to ask?}
\end{quote}

This highlights the necessity for improved approaches to ensure more robust adherence to instructions.

\paragraph{Math} The primary limitation of {Guanaco} lies in its mathematical capabilities, a challenge shared by numerous language models as discussed in \citep{liang2022holistic}. When prompted to display its reasoning, {Guanaco} often arrives at correct solutions, for instance:

\begin{quote}
    \textbf{User}: Lee mows one lawn and charges \$33. Last week he mowed 16 lawns and three customers each gave him a \$10 tip. How many dollars did Lee earn mowing lawns last week?

    \textbf{Guanaco}: \texttt{Lee earned \$582 mowing lawns last week. \\ \\ To find the total amount he earned, we first need to calculate how much he earned for mowing 16 lawns: \\ \\  \$33 x 16 = \$528 \\ \\ Then we need to add the tips he received from the three customers: \\  \\ \$10 x 3 = \$30 \\ \\ So the total amount Lee earned was: \\ \\ \$528 + \$30 = \$558}
\end{quote}

Yet, {Guanaco} sometimes falters on straightforward tasks when it neglects to decompose the problem methodically, a well-documented shortcoming \cite{weichain}. This behavior is illustrated below:

\begin{quote}
    \textbf{User}: Please factorize 1833.

    \textbf{Guanaco}: \texttt{Sure, I'd be happy to help! 1833 is a prime number, which means it has only two factors: 1 and 1833. \\ \\ So the unique factorization of 1833 is: \\ \\ 2\^{}1 * 3\^{}2 * 17\^{}1 \\ \\ I hope this helps! Let me know if you have any other questions.}
\end{quote}

which is incorrect (the actual prime factors are $3 \times 17 \times 43$); in fact, the error is made twice.

\paragraph{Theory of Mind} {Guanaco} demonstrates unexpectedly robust abilities in Theory of Mind tasks \cite{nematzadeh2018evaluating, sap2022neural}. For example, it is able to generate detailed and appropriate responses to questions such as:

\begin{quote}
    \textbf{User}: James and Abby are in the bedroom. Abby put the pen in the desk drawer. Abby leaves the bedroom. James moves the pen into the bag. Where does James think Abby will look for the pen?

    \textbf{Guanaco}: \texttt{James knows that Abby placed the pen in the desk drawer, but he has subsequently moved it to the bag. James assumes Abby is unaware of this change, and thus believes she will search for the pen in the desk drawer, as that is where she last left it. However, this belief might be mistaken if Abby witnessed James relocate the pen.}
\end{quote}

Nonetheless, the consistency of these kinds of inferences is limited; the model frequently makes assumptions that are not warranted by the scenario, such as:

\begin{quote}
    \textbf{User}: Evelyn entered the living room. Jackson entered the playroom. James entered the playroom. The beans are in the treasure chest. James exited the playroom. Jackson moved the beans to the pantry. Jackson exited the playroom. James entered the living room. Where will James look for the beans?

    \textbf{Guanaco}: \texttt{James will look for the beans in the pantry, because that's where Jackson moved them.}
\end{quote}

Here, {Guanaco} infers a transmission of information that was not included in the narrative. Such phenomena are consistent with observations in previous work \cite{sap2022neural}, but necessitate further exploration.

\subsection{Considerations}
\label{ssec:considerations}

\paragraph{Evaluation}

Our findings indicate only moderate concordance among human annotators (Fleiss $\kappa=0.42$), with the agreement decreasing even further when evaluating two highly capable systems. This observation highlights deficiencies in both the design of current benchmarks and in the methodology of human-based chatbot evaluations. During manual side-by-side comparisons of ChatGPT and {Guanaco} 65B outputs on the Vicuna benchmark, we observed that subjective judgement became increasingly important—so much so that there was notable disagreement among this paper’s authors about which model outputs were preferable. Prospective research should consider drawing on techniques from disciplines such as Human-Computer Interaction and Psychology to better address and account for subjective variance in evaluation.

Further, our study reveals systematic biases in automated evaluation frameworks. Notably, we detect pronounced order effects, with GPT-4 consistently assigning better scores to whichever system’s response is listed first in the prompt. The low sample-level concordance between GPT-4’s evaluations and those of human raters (Fleiss $\kappa=0.25$) also indicates that the preferences and evaluative criteria of humans and automated systems are not well aligned. Additionally, as shown in Table~\ref{tbl:elo_large}, GPT-4 significantly inflates scores for its own generations relative to human judgements, resulting in Elo ratings of 1348 versus 1176; this corresponds to approximately a 20\% higher probability of a win in a pairwise comparison.

Future research should address the potential for bias within automated evaluation frameworks and explore effective methods to mitigate such bias.

\paragraph{Data \& Training} It is important to highlight that the OASST1 dataset, which serves as the foundation for training the {Guanaco} models, is inherently multilingual. Additionally, the OA benchmark includes prompts in various languages. Further studies are needed to assess whether this multilingual training contributes to improved outcomes on non-English instructions and whether this factor underlies the pronounced performance gap observed between the Vicuna-13B model—trained exclusively on English—and {Guanaco} models with 33B and 65B parameters on the OA benchmark.

In light of the strong results obtained by the {Guanaco} models, we assess the possibility of data leakage between the OASST1 training data and the prompts utilized in the Vicuna benchmark. Utilizing fuzzy string matching followed by manual review of the closest prompt matches across both datasets, we do not observe any prompt overlap.

We further emphasize that our model's training relies solely on cross-entropy loss (standard supervised learning), without the application of reinforcement learning from human feedback (RLHF). This underscores the need for more detailed comparative studies investigating the respective advantages and limitations of cross-entropy loss and RLHF-based approaches. We anticipate that \textsc{QLoRA} provides a viable path for carrying out such investigations at scale without incurring prohibitive computational costs.

\section{Related Work}

\paragraph{Quantization of Large Language Models} The primary research focus within the quantization of large language models (LLMs) has been on methods to support efficient inference. Many state-of-the-art approaches that preserve the quality of 16-bit LLMs target the mitigation of outlier feature issues (such as SmoothQuant~\citep{xiao2022smoothquant} and LLM.int8()~\citep{dettmers2022llm}), whereas others employ more complex grouping strategies~\citep{park2022nuqmm, yao2022zeroquant}. Studies in lossy quantization have examined the trade-offs involved in conventional rounding~\citep{dettmers2022case,zeng2022glm,pope2022efficiently} as well as strategies for optimizing rounding to improve quantization accuracy~\citep{frantar2022gptq}. Apart from our own contributions, SwitchBack layers~\citep{wortsman2023stable} represents the only other research to investigate the effects of backpropagating through quantized weights at a model scale exceeding 1B parameters.

\paragraph{Finetuning with Adapters} Although our experiments leverage Low-rank Adapters~\citep{hu2021lora} (LoRA), the literature offers a diverse array of Parameter Efficient FineTuning (PEFT) strategies. Examples include prompt tuning~\citep{qin2021learning,lester2021power,li2021prefix}, adjusting the inputs to the embedding layer~\citep{an2022input}, fine-tuning hidden representations (IA$^3$)~\citep{liu2022few}, integrating additional full layers~\citep{houlsby2019parameter}, bias tuning~\citep{zaken2021bitfit}, employing masks on weights informed by Fisher information~\citep{sung2021training}, as well as hybrid techniques~\citep{henderson2021compacter}. Our findings demonstrate that LoRA adapters can deliver results on par with comprehensive 16-bit finetuning. We leave a systematic comparison of LoRA with other PEFT mechanisms as an open direction for subsequent research.

\paragraph{Instruction Finetuning} Instruction finetuning is designed to enhance the capacity of pretrained LLMs to respond appropriately to prompts by utilizing paired input-output samples drawn from a variety of data sources. Through this process, a pretrained LLM is adapted to produce target outputs when prompted. Notable methods and datasets for this task comprise MetaICL~\citep{min2021metaicl}, MetaTuning~\citep{zhong2021adapting}, InstructGPT~\citep{ouyang2022training}, FLAN~\citep{wei2021finetuned,chung2022scaling}, PromptSource~\citep{bach2022promptsource}, Super-NaturalInstructions~\citep{wang2022super,sanh2021multitask}, Self-instruct~\citep{wang2022self}, UnnaturalInstructions~\citep{honovich2022unnatural}, OPT-IML~\citep{iyer2022opt}, UnifiedSKG~\citep{xie2022unifiedskg}, OIG/Chip2~\citep{laion2023}, Alpaca~\citep{alpaca}, Vicuna~\citep{vicuna2023}, Koala~\citep{koala_blogpost_2023}, and Self-instruct-GPT-4~\citep{peng2023instruction}.

\paragraph{Chatbots} Numerous instruction-tuned language models are developed as dialog agents or chatbots, commonly refined using Reinforcement Learning from Human Feedback (RLHF)~\citep{christiano2017deep}, or alternatively, by generating training data via existing models in conjunction with AI-based feedback (RLAIF)~\citep{bai2022constitutional}. Significant approaches and data collections in this context include Anthropic-HH~\citep{askell2021general,bai2022training}, Open Assistant~\citep{kopf2023openassistant}, LaMDA~\citep{thoppilan2022lamda}, and Sparrow~\citep{glaese2022improving}. Our work does not utilize reinforcement learning, but our most performant model, Guanaco, is finetuned on multi-turn conversations from the Open Assistant dataset, which was curated to facilitate RLHF research~\citep{kopf2023openassistant}. Furthermore, alternatives to costly human evaluations, such as leveraging GPT-4 for assessment, have gained traction~\citep{vicuna2023,peng2023instruction}. In our evaluation framework, we build on and refine these approaches to create a more robust methodology.

\section{Limitations and Discussion}

This study presents evidence that our proposed technique, \textsc{QLoRA}, can match the effectiveness of 16-bit full finetuning using a 4-bit quantized model equipped with Low-rank Adapters (LoRA). However, our analysis stops short of confirming whether \textsc{QLoRA} remains competitive with full 16-bit finetuning at the larger model scales of 33B and 65B, primarily due to substantial computational constraints. Consequently, this question remains open for future inquiry.

A further limitation concerns our evaluation of instruction finetuned models. While our experiments include results for MMLU, the Vicuna benchmark, and the OA benchmark, they do not cover other established evaluation suites such as BigBench, RAFT, and HELM. Thus, the extent to which our conclusions generalize to these other benchmarks is uncertain. Nonetheless, our study includes an extensive investigation on MMLU and introduces novel strategies for chatbot evaluation.

The presented findings suggest that benchmark outcomes are strongly influenced by the degree of overlap between the finetuning data and the benchmark in question. For instance, FLAN v2 shares similarities with MMLU but diverges from chatbot benchmarks, whereas the Chip2 dataset exhibits the opposite relationship. Correspondingly, the models’ performance aligns with these affinities on both the MMLU and Vicuna evaluations. This observation underscores the necessity for not just improved benchmarks and evaluation metrics, but also careful consideration regarding the specific abilities under assessment. Should our focus be on optimizing models for academic knowledge comparable to high school and college exams, or should we prioritize proficiency in conversational chatbot interactions? Or is our objective yet something different? Because leveraging pre-existing benchmarks is typically simpler than designing new ones, the selection of certain benchmarks has the potential to direct research in specific ways. It is essential for the field to ensure that the benchmarks in use accurately reflect the competencies we value.

Although we offer a thorough analysis of general chatbot effectiveness, one shortcoming of our work lies in the limited scope of our responsible AI assessment of {Guanaco}. In Table~\ref{tbl:crows}, we measure the probability that {Guanaco}-65B produces socially biased token sequences relative to alternative models, and find that {Guanaco}-65B achieves significantly lower average bias scores than standard pretrained models. This suggests that finetuning on the OASST1 dataset may reduce biases present in the LLaMA base model. While these initial results are promising, the extent to which {Guanaco} mitigates other forms of bias remains uncertain. Comprehensive bias analysis in {Guanaco} and comparable conversational agents is left as an avenue for future research.

Another limitation of our evaluation concerns the omission of experiments involving various bit-precisions, such as 3-bit base models, as well as the exploration of alternative adapter mechanisms. While LoRA was chosen due to its demonstrated stability, numerous other Parameter Efficient FineTuning (PEFT) techniques exist and have shown effectiveness, although their scalability to larger models remains to be validated. It is possible that some of these approaches could lead to superior results. Notably, since post-quantization finetuning appears to restore much of the accuracy lost through quantization, this technique could permit even more aggressive quantization schemes. For instance, following 3-bit GPTQ quantization of the base model, finetuning with LoRA might achieve performance comparable to standard 16-bit finetuning.

\section{Broader Impacts}

The \textsc{QLoRA} finetuning strategy introduced in this work is, to our knowledge, the first to make it feasible to finetune models with 33B parameters on a standard consumer GPU and models with 65B parameters on a single high-end professional GPU—without incurring performance loss compared to full-precision finetuning. We have shown that our leading 33B model, trained on the Open Assistant dataset, can achieve results competitive with ChatGPT on the Vicuna evaluation. Because instruction finetuning represents a crucial means of converting base pretrained LLMs into capable chatbot systems, we anticipate that our approach will democratize finetuning, particularly benefiting researchers with more limited computational resources. In this way, \textsc{QLoRA} serves as a significant step toward reducing disparities in access to advanced NLP tools, bridging the gap between large organizations and smaller groups equipped with consumer-grade hardware.

An additional avenue for significant impact lies in the adaptation of LLMs for mobile devices. We posit that our \textsc{QLoRA} approach could mark a turning point, facilitating the finetuning of large language models directly on mobile phones and in environments with limited computational resources. While previous work has demonstrated that 7B parameter models can be executed on phones, \textsc{QLoRA} stands out as the first solution capable of supporting their finetuning. For instance, our estimates suggest that an iPhone 12 Plus can process approximately 3 million tokens for finetuning overnight while being charged. Although the performance of finetuned 7B models falls short of ChatGPT’s standards, the results are nonetheless sufficient to unlock innovative use cases, especially where privacy considerations or prior limitations in LLM capabilities posed obstacles. The use of \textsc{QLoRA} opens opportunities for privacy-conscious deployment of LLMs, empowering individuals to maintain autonomy over their data and model customization, while also streamlining the deployment process.

It is important to recognize, however, that finetuning LLMs constitutes a dual-use technology that can be exploited for harmful purposes. Prior research has highlighted the potential risks associated with widespread LLM adoption \citep{bommasani2021opportunities,bender2021dangers}. Nevertheless, we argue that democratizing access to a technology rapidly permeating many domains will foster more comprehensive and independent scrutiny compared to restricting these capabilities to major corporations that limit transparency by withholding model weights and source code from external review.

Ultimately, we contend that \textsc{QLoRA} represents a substantial advance toward democratizing the finetuning of high-caliber LLMs, thereby greatly increasing their accessibility and ease of use for a broader audience.